\documentclass[11pt,a4paper]{article}
\usepackage[utf8]{inputenc}
\usepackage[T1]{fontenc}
\usepackage[french]{babel}
\usepackage{amsmath, amssymb, amsthm}
\usepackage{geometry}
\geometry{margin=2.5cm}

\title{Correction du TD 4 : Analyse et Commande des systèmes linéaires}
\author{AU361 - Grenoble INP Esisar}
\date{}

\begin{document}

\maketitle

\section*{Exercice 1 : Synthèse RST pour un système du 1er ordre}

\textbf{1. Modélisation et cahier des charges}\\
La fonction de transfert du processus est $H(p) = \frac{1}{1+p}$. 
On identifie les polynômes du système : $A(p) = p+1$ (degré $n_A=1$) et $B(p) = 1$ (degré $n_B=0$).

Le cahier des charges impose :
\begin{itemize}
    \item \textbf{Rejet asymptotique des perturbations constantes en sortie :} La fonction de sensibilité $S_{yp}(p) = \frac{A(p)S(p)}{A(p)S(p)+B(p)R(p)}$ doit s'annuler pour $p=0$. Puisque $A(0) = 1 \neq 0$, il faut impérativement imposer $S(0)=0$. On fixe donc une partie de $S(p)$ : $S(p) = p \cdot S'(p)$.
    \item \textbf{Pôle de poursuite :} $p_c = -4 \implies P_c(p) = p+4$.
    \item \textbf{Pôles de filtrage :} $p_f = -2 \implies P_f(p) = (p+2)^{n_f}$.
\end{itemize}

\textbf{2. Degré des polynômes}\\
L'équation de Bezout est $A(p)S(p) + B(p)R(p) = P(p)$.\\
Pour avoir un régulateur propre (ou strictement causal), on pose $n_R = n_S$.\\
Le polynôme $S(p)$ s'écrit $S(p) = p$ (pour avoir le moins de degrés possibles, $n_S=1$).\\
On a donc $n_R = 1 \implies R(p) = r_1 p + r_0$.\\
Le polynôme caractéristique $P(p)$ a pour degré $\deg(P) = n_A + n_S = 1+1=2$.\\
Ainsi, le polynôme de filtrage sera de degré 1 : $P_f(p) = p+2$.\\
Le polynôme cible est donc $P(p) = (p+4)(p+2) = p^2 + 6p + 8$.

\textbf{3. Calcul des polynômes $R$, $S$ et $T$}\\
On résout l'équation de Bezout :
$$ (p+1)(p) + 1 \cdot (r_1 p + r_0) = p^2 + 6p + 8 $$
$$ p^2 + p + r_1 p + r_0 = p^2 + (1+r_1)p + r_0 $$
Par identification avec $P(p)$ :
\begin{itemize}
    \item $1+r_1 = 6 \implies r_1 = 5$
    \item $r_0 = 8$
\end{itemize}
Les polynômes du correcteur sont donc :
$$ S(p) = p $$
$$ R(p) = 5p + 8 $$
Pour obtenir un gain statique unitaire sans modifier la dynamique de poursuite, on annule les pôles de filtrage via :
$$ T(p) = \frac{P_c(0)P_f(p)}{B(0)} = 4(p+2) = 4p+8 $$
\textit{(Note : le choix $T(p) = \frac{P(0)}{B(0)} = 8$ est également valide mais ne compense pas le pôle de filtrage dans la poursuite).}

\textbf{4. Transfert en boucle fermée et sensibilité}
\begin{itemize}
    \item \textbf{Transfert BF :} $\frac{Y(p)}{C(p)} = \frac{B(p)T(p)}{P(p)} = \frac{4(p+2)}{(p+4)(p+2)} = \frac{4}{p+4}$.
    \item \textbf{Fonction de sensibilité :} $S_{yp}(p) = \frac{A(p)S(p)}{P(p)} = \frac{p(p+1)}{p^2+6p+8}$.
\end{itemize}

\hrulefill

\section*{Exercice 2 : Structure du correcteur pour un 2nd ordre}

\textbf{1. Modélisation}\\
Le système est $H(p) = \frac{K}{\frac{p^2}{\omega_n^2} + \frac{2\zeta}{\omega_n}p + 1} = \frac{B(p)}{A(p)}$.\\
Le degré de $A(p)$ est $n_A = 2$.

\textbf{2. Condition de rejet de perturbation}\\
On veut rejeter une perturbation de pulsation $\omega = 1$ rad/s. La fonction de sensibilité $S_{yp}(p)$ doit être nulle en $p = \pm j$. Il faut donc que $S(p)$ contienne le facteur $(p-j)(p+j) = p^2+1$.\\
On pose $S(p) = (p^2+1)S'(p)$. Le degré minimal de $S(p)$ est donc augmenté de 2.

\textbf{3. Détermination des degrés}\\
On cherche un correcteur causal ($n_R = n_S$).\\
L'équation de Bezout est de degré $\deg(P) = n_A + n_S = 2 + n_S$.\\
Le nombre d'inconnues est $n_{S'}$ (pour $S'$ unitaire) $+ n_R + 1$.\\
Or $n_S = n_{S'} + 2$, donc $n_R = n_{S'} + 2$.\\
Le nombre total d'inconnues est $n_{S'} + n_{S'} + 2 + 1 = 2n_{S'} + 3$.\\
En égalisant le nombre d'équations et d'inconnues :
$$ \deg(P) = 2 + n_S = 2 + (n_{S'} + 2) = n_{S'} + 4 $$
$$ n_{S'} + 4 = 2n_{S'} + 3 \implies n_{S'} = 1 $$

\textbf{Conclusion sur la structure :}\\
\begin{itemize}
    \item Le polynôme $S(p)$ est de degré $1 + 2 = \textbf{3}$.
    \item Le polynôme $R(p)$ est de degré $\textbf{3}$.
    \item Le polynôme caractéristique $P(p)$ sera de degré $\textbf{5}$.
\end{itemize}

\hrulefill

\section*{Exercice 3 : Moteur asservi en position (Régulateur RST)}

\textbf{1. Modélisation}\\
$H(p) = \frac{K}{p(1+Tp)} = \frac{B(p)}{A(p)}$ avec $n_A = 2$ et $n_B = 0$.

\textbf{2. Synthèse sans erreur statique et rejet perturbation sur mesure}\\
\begin{itemize}
    \item \textbf{Suivi consigne constante :} Implique un gain statique unitaire, obtenu par un choix approprié de $T(p)$ ($T(0) = P(0)/B(0)$).
    \item \textbf{Rejet perturbation constante additive sur la position :} C'est une perturbation de sortie. Il faut $S_{yp}(0) = 0 \implies A(0)S(0) = 0$. Puisque $A(p)$ contient déjà un intégrateur ($A(0)=0$), la condition est naturellement remplie. Aucune partie fixe n'est requise pour $S(p)$.
    \item \textbf{Structure :} On pose $n_R = n_S$. $\deg(P) = 2 + n_S$. Inconnues : $n_S + n_R + 1 = 2n_S + 1$.
    $$ 2 + n_S = 2n_S + 1 \implies n_S = 1 $$
    Les polynômes $R(p)$ et $S(p)$ sont de \textbf{degré 1}, $P(p)$ est de \textbf{degré 3}.
\end{itemize}

\textbf{3. Rejet asymptotique des couples résistants}\\
\begin{itemize}
    \item La perturbation $\Gamma_r$ agit sur l'entrée du système. L'annulation de l'erreur nécessite $B(0)S(0) = 0$.
    \item Puisque $B(0) = K \neq 0$, on doit imposer \textbf{$S(0)=0$}. Le polynôme $S$ doit contenir un intégrateur : $S(p) = p S'(p)$.
    \item \textbf{Conséquence sur les degrés :} $n_S = n_{S'} + 1$. Inconnues : $n_{S'} + n_R + 1$. Pour $n_R = n_S = n_{S'} + 1$, on a $2n_{S'} + 2$ inconnues. Équations : $\deg(P) = 2 + n_S = 3 + n_{S'}$.
    $$ 3 + n_{S'} = 2n_{S'} + 2 \implies n_{S'} = 1 $$
    Les polynômes $R(p)$ et $S(p)$ sont de \textbf{degré 2}, $P(p)$ de \textbf{degré 4}.
\end{itemize}

\hrulefill

\section*{Exercice 4 : Synthèse RST stricte}

\textbf{1. Modélisation et contraintes}\\
$A(p) = p(1+Tp) \implies n_A = 2$. $B(p) = G \implies n_B = 0$.
\begin{itemize}
    \item \textbf{Strictement causal :} $n_R \le n_S - 1$.
    \item \textbf{Rejet $d_1$ (sortie) :} Naturellement rejeté car $A(0)=0$.
    \item \textbf{Rejet $d_2$ (entrée) :} Impose $S(0)=0 \implies S(p) = p S'(p)$.
    \item \textbf{Poursuite :} $P_c(p) = p^2 + 1.4p + 1$ ($\omega_n=1$, $\zeta=0.7$).
    \item \textbf{Filtrage :} Pôles à -1.
\end{itemize}

\textbf{2. Détermination des degrés}\\
On pose $n_R = n_S - 1$. Le nombre d'inconnues est $n_{S'} + n_R + 1 = n_{S'} + (n_S - 1) + 1 = n_{S'} + (n_{S'} + 1 - 1) + 1 = 2n_{S'} + 1$. \\
Le nombre d'équations est $\deg(P) = n_A + n_S = 2 + n_{S'} + 1 = n_{S'} + 3$.
$$ 2n_{S'} + 1 = n_{S'} + 3 \implies n_{S'} = 2 $$
Si $n_{S'} = 2$, alors $n_S = 3$ et $n_R = 2$.
La condition de stricte causalité est bien respectée. 
\begin{itemize}
    \item $\deg(S) = \textbf{3}$
    \item $\deg(R) = \textbf{2}$
    \item $\deg(P) = \textbf{5}$
\end{itemize}
Ainsi $P_f(p) = (p+1)^3$. Les polynômes se déduisent par identification de Bezout : $A(p)S(p) + B(p)R(p) = P_c(p)P_f(p)$.

\hrulefill

\section*{Exercice 5 : Régulation de température d'une enceinte thermique}

\textbf{1. Modèle isolé ($\alpha_2 = 0$)}\\
La mise en équation donne la fonction de transfert entre la puissance et la température :
$$ H(p) = \frac{\alpha_3}{p(p + \alpha_1 + \alpha_3)} = \frac{\frac{\alpha_3}{\alpha_1+\alpha_3}}{p\left(1 + \frac{p}{\alpha_1+\alpha_3}\right)} $$
On identifie $G = \frac{\alpha_3}{\alpha_1+\alpha_3}$ et $T = \frac{1}{\alpha_1+\alpha_3}$. L'erreur statique en boucle fermée pour un échelon est nulle car $H(p)$ possède un intégrateur.

\textbf{2. Système avec pertes (Approximation)}\\
Avec $\alpha_1=0.9$, $\alpha_2=0.01$, $\alpha_3=0.1$, le transfert exact possède 2 pôles ($p_1 \approx -0.001$ et $p_2 \approx -1.01$). Le pôle lent domine, permettant une approximation du 1er ordre.
La constante de temps est $\tau = 1000$ s et le gain statique est $100$.
$$ H_{approx}(p) = \frac{100}{1 + 1000p} $$
Ici $A(p) = 1+1000p$ et $B(p) = 100$.

\textbf{3. Synthèse RST}\\
\begin{itemize}
    \item \textbf{Temps de réponse 5\% de 900s :} Pour un 1er ordre en BF, $t_{r5\%} \approx 3\tau_{bf} \implies \tau_{bf} = 300$ s. Pôle de poursuite : $p_c = -1/300$.
    \item \textbf{Rejet perturbation constante :} Perturbation de sortie (température externe) $\implies S(p)$ doit forcer $S_{yp}(0)=0$. Comme $A(0)=1$, il faut $S(0)=0 \implies S(p) = p(p+s_1)$. (Degré $n_S = 2$).
    \item \textbf{Tous pôles égaux :} $P(p)$ de degré 3, donc $P(p) = \left(p + \frac{1}{300}\right)^3$.
\end{itemize}

\textbf{4. Calcul de Bezout :}\\
En utilisant le modèle normalisé $H(p) = \frac{0.1}{p+0.001}$, l'identification polynomiale $(p+0.001)(p^2+s_1p) + 0.1(r_1p+r_0) = \left(p+\frac{1}{300}\right)^3$ donne par développement :
\begin{itemize}
    \item $s_1 = 0.009$
    \item $r_1 \approx 2.43 \times 10^{-4}$
    \item $r_0 \approx 3.7 \times 10^{-7}$
\end{itemize}

\textbf{5. Fonction de Sensibilité :}\\
Le diagramme de Bode asymptotique de $S_{yp}(p) = \frac{p(p+0.001)(p+0.009)}{\left(p+\frac{1}{300}\right)^3}$ :
\begin{itemize}
    \item Débute en basses fréquences avec une pente de +20 dB/décade (effet du zéro à l'origine).
    \item Subit deux cassures vers le haut (+20 dB/déc chacune) aux pulsations $\omega = 0.001$ et $\omega = 0.009$.
    \item Subit une triple cassure vers le bas (-60 dB/déc) au pôle dominant $\omega = 0.0033$ ($1/300$).
    \item Finit sur un gain asymptotique de 0 dB en hautes fréquences.
\end{itemize}

\end{document}